\section{Recovering the character labels}
\label{sec:blind}

Suppose the character list is itself unknown. A coupling of larger size can contain a smaller coupling as a direct summand, inheriting every one of its zeros. At equal or smaller size, even a one-sided comparison determines all the labels. We allow repetitions and include the primitive character of conductor one, with completion $\xi(s)=\tfrac12 s(s-1)\pi^{-s/2}\Gamma(s/2)\zeta(s)$.

\begin{theorem}[Blind recovery from a one-sided tail]\label{thm:blind}
Assume the Dirichlet critical-line input for the characters in two lists $\mathcal A,\mathcal B$ of sizes $n_A,n_B$. Let $A,B$ be Hermitian couplings on those lists. Suppose that all but finitely many zero multiplicities of $P_A$ occur among those of $P_B$. Then $n_B\ge n_A$. If $n_B=n_A$, the primitive character multisets coincide. After relabelling, the two determinant polynomials in the distinct label variables coincide, including all their grouped principal-minor coefficients. In particular, distinct labels determine the individual principal minors, and a nonnegatively weighted graph Laplacian is determined up to simultaneous relabelling when the labels are distinct.
\end{theorem}

\begin{proof}
The real-zero count of a Hermitian coupling on an interval with regular endpoints differs from the sum of the input counts by at most its size, by the inertia identity for the coupled determinant. The two-sided Riemann--von Mangoldt count for each fixed primitive completion on $[-T,T]$ is $(T/\pi)\log T+O_{\chi}(T)$. Passing to arbitrarily large regular endpoints gives $n_B\ge n_A$. Write $n=n_A=n_B$ for the rest of the proof.

Put $s=\tfrac12+iz$ and $C_A(s)=P_A((s-\tfrac12)/i)$, and define $C_B$ in the same way. For a list $\mathcal L$ let $Q_{\mathcal L}$ be the product of its conductors, $a_{\mathcal L}$ the sum of its parities (each $0$ or $1$), and $h_{\mathcal L}$ the number of copies of the conductor-one label. Define
\[
 c=\tfrac12\log(Q_{\mathcal B}/Q_{\mathcal A}),\qquad
 \beta=\tfrac12(a_{\mathcal B}-a_{\mathcal A})+2(h_{\mathcal B}-h_{\mathcal A}),\qquad
 G(s)=\frac{\prod_{\chi\in\mathcal B} L(s,\chi)}{\prod_{\chi\in\mathcal A} L(s,\chi)}.
\]
Every zero of $C_A$ is on the critical line. The one-sided hypothesis therefore supplies a monic polynomial $W$ cancelling all the poles of $C_B/C_A$ with their orders. Write $w=\deg W$, $e=w+\beta$ and $E=W C_B/C_A$. This entire function has order at most one. Indeed, both coupled functions are finite sums of products of entire functions of order at most one and their derivatives. In their genus-one Hadamard representations the denominator zeros cancel against those of $W C_B$, leaving a linear exponential and a canonical product over a submultiset of the zeros of $W C_B$. The counting bound $O_\epsilon(r^{1+\epsilon})$ on this submultiset gives order at most one.

Uniformly for $\Re s\ge 2$ and $|s|\to\infty$, there is a fixed $d\ne0$ for which
\begin{equation}\label{eq:blind-asymptotic}
 E(s)=d\,e^{cs}s^e G(s)\bigl(1+O(1/\log |s|)\bigr),
\end{equation}
using the principal logarithm in this half-plane. To verify the formula, divide each determinant by the product of its completions. The resulting determinant has diagonal logarithmic derivatives $\tfrac i2\log(s/2)+O(1)$ plus the fixed matrix. Here $L'/L$ is bounded by absolute convergence, the digamma estimate is uniform, and the extra logarithmic derivatives $1/s+1/(s-1)$ for $\xi$ are bounded. The two determinants both have leading term $(\tfrac i2\log(s/2))^n$ with relative error $O(1/\log|s|)$. In the ratio of the completion products, the equal numbers of gamma factors cancel their common part, and $\Gamma((s+1)/2)/\Gamma(s/2)=(s/2)^{1/2}(1+O(1/|s|))$. The additional $s(s-1)$ at each conductor-one occurrence contributes two powers. The conductor exponentials supply $e^{cs}$, the monic polynomial contributes $s^w(1+O(1/|s|))$, and all nonzero constant factors are absorbed into $d$.

Write $\kappa=|c|$. Since $G$ is bounded on $\Re s\ge 2$, the right half-plane satisfies $|E(s)|\ll (1+|s|)^e e^{c\Re s}$, enlarging the constant on a compact set. Phase normalization makes $P_A$ and $P_B$ real entire in $z$, hence $|C_j(1-\overline s)|=|C_j(s)|$. Reflection gives $|E(s)|\ll (1+|s|)^e e^{c(1-\Re s)}$ on $\Re s\le -1$, since the modulus of a monic polynomial of degree $w$ is comparable to $(1+|s|)^w$ off a compact set. In the intervening strip consider $E(s)/(s+3)^e$ with the analytic branch positive on the positive real axis. This quotient is bounded on both vertical boundaries. For $0<\alpha<\pi/3$, apply the maximum principle on expanding rectangles to the quotient multiplied by $\exp[-\delta\cos(\alpha(s-1/2))]$. The real part of the cosine is at least $\cos(3\alpha/2)\cosh(\alpha\Im s)$ throughout the strip, and the finite order bound controls the horizontal sides. Letting $\delta\downarrow0$ and combining the three regions yields
\begin{equation}\label{eq:blind-growth}
 |E(s)|\le C(1+|s|)^e\exp(\kappa|\Re s|)\qquad(s\in\C).
\end{equation}
Retaining the signed power $e$ in this estimate will remove the polynomial growth when the function is translated.

Choose $T_j\to+\infty$ so that $p^{-iT_j}\to1$ for every prime $p$ and $e^{icT_j}\to1$. For any finite list of these phases, simultaneous recurrence follows from compactness of the torus (or simultaneous Dirichlet approximation). Taking successively longer finite lists gives the sequence. Absolute convergence gives $G(s+iT_j)\to G(s)$ locally uniformly for $\Re s>1$, since the Dirichlet coefficients of the quotient are absolutely summable with weights $m^{-\sigma}$ for every $\sigma>1$. By \eqref{eq:blind-growth}, the entire functions $T_j^{-e}E(s+iT_j)$ are locally bounded. A subsequence converges locally uniformly to an entire function $H$ with
\[
 |H(s)|\le C e^{\kappa |\Re s|},\qquad
 H(s)=d_0 e^{cs}G(s)\quad(\Re s\ge2),\qquad d_0=d\exp(i\pi e/2)\ne0.
\]
Indeed, $(1+|s+iT_j|)/T_j\to1$ for each fixed $s$, and the second assertion follows from \eqref{eq:blind-asymptotic}.

Write $G(s)=\sum_{m\ge1} b_m m^{-s}$. For fixed $m$ put $\lambda=c-\log m$. On $\Re s=2$ the vertical mean of $H(s)e^{-\lambda s}$ equals $d_0 b_m$ by absolute convergence and orthogonality of distinct exponential frequencies. The same mean on $\Re s=-u$ has the same value for every $u>0$. To see this, integrate around the rectangle with those two vertical sides and horizontal sides at $\pm iY$, and divide by $2iY$. The horizontal integrals are $O_u(1)$ because of the global vertical-line bound for $H$. Consequently
\[
 |d_0 b_m|\le C\exp((\kappa+\lambda)u).
\]
Letting $u\to\infty$ proves that $G$ is a Dirichlet polynomial supported on $m\le e^{c+|c|}=\max(1,Q_{\mathcal B}/Q_{\mathcal A})$.

In particular $b_p=0$ for all sufficiently large primes. On the other hand, the prime coefficient of the quotient is
\[
 b_p=\sum_{\chi\in\mathcal B}\chi(p)-\sum_{\chi\in\mathcal A}\chi(p).
\]
For every unit residue class modulo the least common multiple of the conductors, Dirichlet's theorem supplies arbitrarily large primes in the class. Thus the two character sums agree on the full unit group. Distinct primitive characters induce distinct characters of this group, so orthogonality identifies all multiplicities. After the labels have been matched, the fixed-list one-sided rigidity theorem gives equality of the grouped determinant polynomials. This proves the assertion.
\end{proof}

\begin{corollary}[Finite corruption and unknown size]\label{cor:blind-symmetric}
A coupled zero divisor with finitely many arbitrary insertions and deletions determines its size, its multiset of primitive labels, and its grouped determinant polynomial. For one-sided containment, the size condition is sharp. In every larger size one can preserve the entire smaller divisor and introduce extra labels by taking a block-diagonal coupling with the original matrix as one block.
\end{corollary}
\begin{proof}
Apply the size inequality in both directions, then the equal-size conclusion. The block-diagonal determinant and the product of its vertex completions both factor, giving the stated example.
\end{proof}

\begin{corollary}[Finite windows for a bounded catalogue]\label{cor:blind-finite}
Fix a finite catalogue of primitive characters, a maximum list size $N$, a norm bound $M$, and an integer erasure budget $K\ge0$. There is $T<\infty$ such that the following implication holds for all lists drawn from the catalogue with $n_B\le n_A\le N$ and matrices of norm at most $M$. If at most $K$ zero multiplicities of $P_A$ in $[-T,T]$ are missing from $P_B$, then the lists coincide as multisets and the grouped determinant polynomials agree under relabelling.
\end{corollary}
\begin{proof}
First exclude unequal label multisets. If no window excluded them uniformly, choose counterexamples on $[-j,j]$. Finitely many choices of lists and compactness of the two matrix balls give a subsequence with fixed distinct multisets and convergent matrices. Both limiting determinants are nonzero and all their zeros are real. Around any finite set of distinct limit zeros of the first determinant choose disjoint disks whose closures contain no other zeros of either limit. Choose boundaries free of zeros of both functions. Rouch\'e's theorem preserves the two individual zero counts in each disk along the subsequence. The sum of the positive parts of their differences is at most $K$, since the approximating zeros are real, lie in the observed window in these disks, and the actual missing count dominates this sum. Exhausting the limit zeros gives one-sided tail discrepancy at most $K$. Theorem~\ref{thm:blind} now contradicts the choice of lists. Once the multisets agree, use the fixed-list finite-window theorem, taking the maximum of its thresholds over the finitely many lists. A finite-search interpretation of a sufficient window is given below.
\end{proof}
